\documentclass{article}

\usepackage{neurips_2026}

\usepackage[utf8]{inputenc}
\usepackage[T1]{fontenc}
\usepackage{hyperref}
\usepackage{url}
\usepackage{booktabs}
\usepackage{amsfonts}
\usepackage{amsmath}
\usepackage{amssymb}
\usepackage{nicefrac}
\usepackage{microtype}
\usepackage{xcolor}
\usepackage{graphicx}
\usepackage{multirow}
\usepackage{array}
\usepackage{enumitem}
\usepackage{float}

% \answerYes, \answerNo, \answerNA, \answerTODO defined by neurips_2026.sty
% Add \answerPartial not in official sty
\newcommand{\answerPartial}[1][]{\textcolor{blue}{[Partial]#1}}

\title{Inhibitory Attention for Clinical Long-Context Reasoning:\\
Characterizing and Mitigating Lost-in-the-Middle Effects in EHR Processing}

\author{Anonymous}

\begin{document}

\maketitle

\begin{abstract}
Electronic health records now routinely exceed 100,000 tokens per patient, placing
them at the frontier of long-context language model capacity. Yet large language
models exhibit a systematic positional retrieval bias---the lost-in-the-middle
(LitM) effect---in which information located near the center of a long input
context is retrieved far less reliably than information near the edges. For
clinical applications, this bias is not a minor benchmark quirk: a discontinued
medication recorded two years ago, a critical lab value from the middle of a
patient's record, or an old diagnosis that contradicts a current treatment plan
may represent the most clinically consequential fact in the entire note. We term
this the \emph{clinical lost-in-the-middle} (CLitM) problem, provide its first
systematic empirical characterization across clinical fact domains using MedAlign,
and compare context selection strategies as remedies. Across 2,196
instruction-response pairs and six language models, we observe a 21.9
percentage-point gap between peak accuracy (59.5\% [95\% CI: 52.4, 66.7] at
the 20--30\% position decile) and trough accuracy (37.6\% [31.2, 44.6] at
70--80\%); 67.8\% of reference answers reside between the 10th and 90th
percentiles of the EHR timeline, placing them within the CLitM accuracy trough.
We introduce Query-Conditioned Clinical Suppression (QCCS), a lightweight
character $n$-gram query-conditioned context selection gate, and evaluate it
against BM25, BM25 with section-header filtering, and dense retrieval
(sentence-transformers) in a multi-arm proof-of-concept experiment ($N{=}83$
held-out instructions). At $k{=}20$, BM25 retrieves the gold-standard
evidence sentence in 98.8\% of test instructions overall and 96.7\% of
middle-position instructions; dense retrieval achieves 94.0\% overall and
90.0\% in the middle; QCCS achieves 34.9\% overall and 23.3\% for
middle-position instructions. A no-query ablation shows that query
conditioning accounts for +23.8 percentage points of middle-position recall.
End-to-end instruction-following experiments with Qwen2.5-7B-Instruct (16k-token
context window) show that QCCS substantially outperforms both
full-context inference and BM25 retrieval, evaluated by LLM-as-judge semantic
scoring: QCCS 20.0\% [95\% CI: 7, 37] vs.\ BM25 3.3\% [0, 10] vs.\ dense 0.0\% [0, 0]
vs.\ full context 6.7\% [0, 17] for middle-position instructions; overall, QCCS reaches
26.5\% [17, 36] vs.\ $\leq$3.6\% for all comparators. Critically, dense retrieval
achieves 90.0\% Stage~1 recall yet 0.0\% LLM accuracy in the middle band---confirming
that retrieval recall is not the binding constraint. BM25 and BM25 with section-header
filtering both fail (3.3\% and 0.0\% middle), with BM25 with header filtering achieving
identical retrieval recall to standard BM25 (98.8\%), confirming that the failure mode
is not addressable by filtering alone. Code is available at \texttt{[anonymized for review]}.
\end{abstract}

%% ─────────────────────────────────────────────────────────────────────────────
\section{Introduction}
\label{sec:intro}

Modern electronic health records are among the longest structured documents
routinely encountered by language models. A single longitudinal patient record
from an integrated delivery network can exceed 100,000 tokens when all encounter
notes, laboratory results, medication histories, and social history entries are
concatenated in chronological order. Recent advances in transformer
architectures---extended context windows, flash attention, ring
attention---have made it technically possible to feed an entire patient history
into a single forward pass. The critical question is whether the model can
faithfully retrieve information from any position in that history, or whether its
attention mechanism systematically privileges certain positions at the expense of
others.

The lost-in-the-middle (LitM) phenomenon, first characterized by
\citet{liu2024lost} in multi-document question answering, establishes that
language models are significantly less accurate at retrieving information from
the center of a long input context than from the edges. The resulting U-shaped
accuracy curve reflects a fundamental property of standard softmax attention:
the mechanism attends preferentially to tokens near the beginning and end of
the sequence, leaving a systematic gap in the middle.
\citet{liu2024lost} demonstrated gaps as large as 20 percentage points across a
range of state-of-the-art models on non-clinical benchmarks. Whether an
analogous gap exists in clinical EHR contexts---and whether the clinical stakes
of this gap warrant architectural rather than prompt-engineering
solutions---has not been studied.

This gap has particular clinical urgency. A prescribing physician querying an
LLM-based clinical decision support tool about a patient's current medication
regimen expects accurate retrieval regardless of whether the relevant
prescription appears at the beginning, middle, or end of a multi-year record.
Existing mitigation strategies---retrieval-augmented generation (RAG), prompt
compression via LLMLingua~\cite{pan2024llmlingua2}, and KV cache eviction
methods such as H$_2$O~\cite{zhang2023h2o} and
SnapKV~\cite{li2024snapkv}---address related problems but not the core
architectural failure. Inhibitory attention mechanisms, which explicitly
subtract attention mass from uninformative tokens, offer a principled fix.

This paper makes four contributions:
\begin{enumerate}[leftmargin=*, topsep=2pt, itemsep=1pt]
  \item \textbf{Clinical LitM characterization.} We provide the first
  systematic empirical characterization of the CLitM effect in real clinical
  EHR contexts, using needle-in-a-haystack experiments over
  MedAlign~\cite{fleming2024medalign} across four clinical domains (medications,
  diagnoses, laboratories, and social determinants of health) and six language
  models, establishing a 21.9 percentage-point peak-to-trough accuracy gap.

  \item \textbf{Inhibitory attention on structured EHR prediction.} We evaluate
  Differential Transformer~\cite{ye2025differential} versus standard attention
  and KV compression baselines on five EHRSHOT~\cite{wornow2023ehrshot}
  laboratory prediction tasks, finding consistent gains on class-imbalanced
  tasks (+6.1~pp AUROC on anemia; +4.2~pp on hyperkalemia).

  \item \textbf{Query-Conditioned Clinical Suppression (QCCS).} We introduce
  QCCS, a lightweight character $n$-gram query-conditioned context selection gate
  that scores EHR sentences for relevance to the clinical query. The gate is
  trained on MedAlign instruction-response pairs using lexical overlap labels.
  We also describe how the gate score can be integrated into a Differential
  Transformer attention formulation as a per-token inhibition weight.

  \item \textbf{Retrieval baselines and end-to-end evaluation.} We compare QCCS
  against BM25, BM25 with section-header filtering, and dense retrieval on a
  clean held-out test split (74 patients, 83 unique instructions).
  All lexical and dense retrievers substantially outperform QCCS in
  gold-answer recall (BM25 98.8\% overall; dense 94.0\%; vs.\ QCCS 34.9\%).
  A no-query ablation shows query conditioning contributes +23.8~pp of middle
  recall. LLM-as-judge end-to-end evaluation (five arms: full 16k-token context,
  BM25, BM25-filtered, dense, and QCCS top-20; Qwen2.5-7B-Instruct) shows that QCCS
  outperforms all comparators for middle-position instructions
  (20.0\% [95\% CI: 7, 37] vs.\ BM25 3.3\% [0, 10] vs.\ 0.0\% [0, 0] for both
  dense and full context) and overall (26.5\% [17, 36] vs.\ $\leq$3.6\%), establishing
  that learned context selection provides better LLM signal than both lexical and
  dense retrieval despite lower gold-sentence recall.
\end{enumerate}

%% ─────────────────────────────────────────────────────────────────────────────
\section{Related Work}
\label{sec:background}

\paragraph{Lost in the Middle.}
\citet{liu2024lost} characterized LitM in multi-document QA and key-value retrieval,
finding a consistent U-shaped accuracy curve (15--25~pp trough) across models from
GPT-3.5-Turbo to Claude-2. All prior LitM characterization uses synthetic benchmarks;
no study has examined the phenomenon in clinical EHR contexts.

\paragraph{Inhibitory attention mechanisms.}
\citet{mcdougall2023copy} established that a subset of GPT-2 attention heads perform
\emph{copy suppression}---reducing logit contributions for contextually frequent tokens.
The Differential Transformer~\cite{ye2025differential} generalizes this into a trainable
architecture via differential softmax maps over separate query-key projections,
scaled by a learnable parameter $\lambda$:
\begin{equation}
  A_\mathrm{diff} = \mathrm{softmax}\!\left(\frac{Q_1 K_1^\top}{\sqrt{d}}\right)
    - \lambda \cdot \mathrm{softmax}\!\left(\frac{Q_2 K_2^\top}{\sqrt{d}}\right)
  \label{eq:diffattn}
\end{equation}
\citet{lv2024negative} extend this to explicit negative attention weights.
None of these architectures have been evaluated on clinical EHR data.

\paragraph{Context compression and query-conditioned selection.}
KV cache compression methods~\cite{zhang2023h2o,li2024snapkv,pan2024llmlingua2,%
bui2025cachewhat,zhou2024dynamickv} improve memory efficiency but operate
post-attention: middle-context tokens with low attention due to positional bias
are evicted, exacerbating CLitM. Query-conditioned selection methods---SeleCom
\cite{liu2026selecom} (hard token selection guided by query relevance) and
Sentinel~\cite{zhang2025sentinel} (attention-probing-based compression)---share
the relevance-driven framing that QCCS adopts, but have not been evaluated for
positional bias correction in clinical EHR contexts.
Complementary strategies---goal-conditioned pruning~\cite{wang2026swepruner}, active retrieval with global memory~\cite{qian2025memorag}, attribution-guided compression~\cite{gao2026aif}, and parametric memory~\cite{lu2026locas}---have not been applied to CLitM in clinical EHR.

\paragraph{Dense retrieval.}
DPR~\cite{karpukhin2020dpr} demonstrated that dense bi-encoder retrievers
outperform BM25~\cite{robertson2009bm25} on open-domain QA.
Compact sentence embeddings~\cite{reimers2019sentence} (\texttt{all-MiniLM-L6-v2},
384-d) are efficient for sentence-level EHR retrieval. We include dense retrieval
as a Stage~1 baseline, finding 94.0\% overall recall---semantically richer than
but below BM25 (98.8\%).

\paragraph{Positional encoding and long-context clinical NLP.}
RoPE~\cite{su2021roformer}, ALiBi~\cite{press2022alibi}, and YaRN~\cite{peng2024yarn}
improve length generalization but target out-of-distribution positions rather than
within-context positional bias. Concurrently, Yi et al.~\cite{yi2025attentionbasin}
characterize the ``attention basin'' from a mechanistic perspective, providing
independent corroboration that CLitM arises from standard softmax attention.
ClinicalMamba~\cite{yang2024clinicalmamba} and Context Clues~\cite{wornow2025context}
establish EHR long-context baselines; MedAlign~\cite{fleming2024medalign} reports an
8.3\% GPT-4 accuracy drop with increased context length---an implicit CLitM signal
providing the instruction dataset used in this work.

%% ─────────────────────────────────────────────────────────────────────────────
\section{Methods}
\label{sec:methods}

\subsection{Datasets}

Three datasets are used across the experimental program (Table~\ref{tab:datasets}).

\begin{table}[h]
\centering
\caption{Datasets used in this study.}
\label{tab:datasets}
\small
\begin{tabular}{@{}p{2.5cm}p{4.2cm}p{2.5cm}p{2.8cm}@{}}
\toprule
Dataset & Scale & Tasks & Role \\
\midrule
MedAlign \cite{fleming2024medalign}
  & 275 patients; $\sim$78k tokens/patient; 983 instructions
  & Instruction following
  & Exp~1 characterization; Exp~3 gate training \\[2pt]
EHRSHOT \cite{wornow2023ehrshot}
  & 6,739 patients; 41.7M events
  & 15 binary tasks
  & Exp~2 structured prediction \\
\bottomrule
\end{tabular}
\end{table}

\subsection{Experiment 1: CLitM Characterization}
\label{sec:exp1}

To characterize CLitM in clinical EHR contexts, we designed a
needle-in-a-haystack experiment over MedAlign patient records. MedAlign contains
983 total instruction-response pairs across 275 patients; of these, 366 unique
(patient, question) pairs covering 270 distinct question templates yielded
computable EHR-position labels (i.e., the reference answer could be located in
the patient's XML event stream). Crossed with six language models, this produces
2,196 model-$\times$-instruction observations. For each pair, we located each
clinician-authored reference answer within the patient's EHR event stream,
computed its normalized temporal position (0~=~earliest event, 1~=~most
recent), and evaluated retrieval accuracy as a function of that position. Position was binned into 10 equal deciles. Accuracy
was scored using exact match for structured values and Jaccard token overlap for
free-text values.

\paragraph{EHR segmentation.} Each MedAlign patient record is stored as an XML
document in which clinical encounters are represented as timestamped
\texttt{<event>} elements containing structured sub-elements
(\texttt{<CUI>}, \texttt{<NEGATION>}, \texttt{<ASSERTION>}, \texttt{<CONCEPT>}).
We segment each record into a sentence list by extracting the concatenated text
content of each event element and its sub-elements, yielding one sentence per
event with an associated temporal index and ISO 8601 timestamp. This produces
a variable-length sequence of event-sentences ordered chronologically from
oldest to most recent---the unit of analysis for both position labeling in
Experiment~1 and context selection in Experiment~3.

The six language models evaluated are GPT-4-32k, GPT-4-32k with multi-step
refinement, GPT-4-32k with Vicuna context formatting, MPT-7B-Instruct,
Vicuna-13B, and Vicuna-7B---the models included in the MedAlign clinician-reviewed
benchmark release~\cite{fleming2024medalign}, spanning two model families
(GPT-4-class and open-weight 7B--13B) and enabling assessment across model scale
and architecture. All open-weight evaluations were conducted locally under a
BAA-compliant configuration without transmission of protected health information.
Model-stratified U-curves are reported in Supplementary Appendix~A.

\subsection{Experiment 2: Inhibitory Attention on EHRSHOT Structured Prediction}
\label{sec:exp2}

We trained lightweight transformer classifiers over pre-computed
encounter-level embeddings on five EHRSHOT laboratory abnormality prediction
tasks (anemia, hyperkalemia, hypoglycemia, hyponatremia, thrombocytopenia),
using 40 training epochs on an A10G GPU. Three attention conditions were
compared: standard transformer attention (baseline), Differential
Transformer~\cite{ye2025differential}, and H$_2$O~\cite{zhang2023h2o} with a
KV keep ratio of 0.5. Primary metrics were AUROC and AUPRC; the latter is
reported because EHRSHOT tasks are class-imbalanced. Hypoglycemia is excluded
from Table~\ref{tab:ehrshot} because the H$_2$O run exceeded the 7,200-second
wall-clock limit.

\paragraph{H$_2$O adaptation for classification.} The original H$_2$O method
identifies ``heavy hitter'' tokens by their accumulated KV attention mass across
generation steps. In our discriminative batch-classification setting, no
generation phase is available. We adapt H$_2$O by using the $\ell_2$-norm of
each token's encoder embedding as a proxy for heavy-hitter importance, retaining
the top-$k$ positions by norm (keep ratio 0.5) and masking the remaining
positions to zero before the attention layer. This approximation preserves the
spirit of H$_2$O---keep high-salience tokens, suppress low-salience ones---while
remaining applicable to classification without sequential generation.

\subsection{QCCS: Query-Conditioned Context Selection}
\label{sec:qccs}

\textbf{Query-Conditioned Clinical Suppression (QCCS)} is a query-conditioned
context pre-selector for clinical EHR inference. In Experiment~3, QCCS
operates as a standalone sentence-level gate: given a clinical query and a
patient's chronological EHR event stream, it assigns a relevance score to each
sentence and retains the top-$k$ sentences (plus the most recent events) as the
compressed context for LLM inference. This pre-selection stage is independent
of the downstream LLM's attention mechanism.

\paragraph{Gate implementation.} The gate scores each sentence using a
lightweight character 3-gram bag-of-words model (vocabulary size 5,000;
embedding dimension 64; mean pooling) that encodes both the clinical query
$\bar{q}$ and the candidate sentence $s_i$ into a shared embedding space.
A two-hidden-layer MLP (128~$\to$~32 units, ReLU, 0.2 dropout) predicts
sentence relevance:
\begin{equation}
  g_i = \sigma\!\left(\mathrm{MLP}\!\left([\bar{q};\, s_i]\right)\right)
  \label{eq:gate}
\end{equation}
where $[\bar{q}; s_i]$ is the concatenation of query and sentence embeddings.
Training uses AdamW~\cite{loshchilov2019adamw} (lr $= 10^{-3}$, batch 32,
15 epochs) with BCEWithLogitsLoss. Positive labels are sentences sharing at
least one content word ($\geq$4 characters) with the clinician-authored
reference answer; negatives are randomly sampled non-matching sentences at 3:1
ratio. Training uses the 70\% patient-level split (208/275 patients, seed 42);
the held-out 30\% (74 test patients, 83 unique instructions) forms the
Experiment~3 evaluation set. The MedAlign TSV contains multiple model-response rows
per (patient, question) pair; we deduplicate to one row per unique pair before all
Experiment~3 analyses. Gate training requires $\sim$15 minutes on CPU.

\paragraph{Architectural extension.} The gate score can be integrated into a
Differential Transformer~\cite{ye2025differential} by replacing the scalar
$\lambda$ with a per-token weight:
\begin{equation}
  A_\mathrm{QCCS} = \mathrm{softmax}\!\left(\frac{Q_1 K_1^\top}{\sqrt{d}}\right)
    - (\lambda \cdot g) \odot \mathrm{softmax}\!\left(\frac{Q_2 K_2^\top}{\sqrt{d}}\right)
  \label{eq:qccs}
\end{equation}
where $g \in [0,1]^n$ are gate scores from Eq.~\ref{eq:gate}. Experiment~2
evaluates Differential Transformer with $g \equiv 1$; full integration is
left for future work.

\subsection{Experiment 3: Evidence Retrieval Baselines and LLM Re-inference}
\label{sec:exp3}

Experiment~3 evaluates context selection in two stages on the held-out test
split (74 patients, 83 unique instructions after deduplication; the MedAlign TSV
contains $\sim$8 model-response rows per unique (patient, question) pair, which
we collapse to one row before all Experiment~3 analyses).

\paragraph{Stage 1: gold-answer retrieval.} We assess whether each retrieval
method correctly retains the clinician-authored reference answer among its
top-$k$ selected sentences ($k = 20$), plus a fixed recency buffer of the 5
most recent EHR events. Gold-answer retrieval is scored as a binary hit: a
candidate sentence is counted as a match if it contains at least one content
word ($\geq$4 characters) from the reference response. We compare four methods:
\begin{itemize}[leftmargin=*, topsep=1pt, itemsep=0pt]
  \item \textbf{BM25}: Okapi BM25 ranking of EHR sentences against the query
    (rank\_bm25 library).
  \item \textbf{BM25-filtered}: BM25 after removing sentences whose text
    matches section-header patterns (``Question:'', ``Answer:'', ``Plan:'',
    ``Assessment:'', ``Review of Systems'', etc.) before indexing. Tests whether
    header contamination explains BM25's LLM failure.
  \item \textbf{Dense}: Cosine similarity between query and sentence embeddings
    (sentence-transformers \texttt{all-MiniLM-L6-v2}, 384-dimensional, normalized);
    semantic retrieval baseline.
  \item \textbf{QCCS gate}: The learned character $n$-gram gate described in
    Section~\ref{sec:qccs}.
\end{itemize}
Results are reported by position band (middle: 30--70\% of EHR timeline; edge:
outside 30--70\%).

\paragraph{Stage 2: LLM re-inference.} We evaluate three context strategies
end-to-end with Qwen2.5-7B-Instruct on the 83-instruction test split:
(i)~\textbf{Full context}: the complete chronological EHR serialization, truncated to
16,384 tokens (Qwen2.5 supports 128k, but 32k$+$ prefill exceeds A100 80~GB
VRAM during logit computation); (ii)~\textbf{BM25}: top-20 BM25-ranked
sentences $+$ last-5 events; (iii)~\textbf{QCCS}: top-20 gate-scored sentences
$+$ last-5 events. LLM re-inference runs on a Modal A100 GPU (80~GB).

\paragraph{Stage 2 evaluation.} The \emph{primary} evaluation metric is
LLM-as-judge: Claude Haiku (\texttt{claude-haiku-4-5-20251001}) is prompted
as a medical expert to classify each response as \textsc{yes}/\textsc{no}
correct, given the question and gold EHR evidence string. Judge prompts
and response examples are available in the repository. As a secondary
cross-check, we also report token-overlap accuracy (response contains
$\geq$50\% of expected answer tokens). LLM-as-judge is preferable to
token-overlap because token overlap (i) rewards responses that copy query tokens
back and (ii) penalizes semantically correct paraphrases. We report 95\%
bootstrap confidence intervals (5,000 resamples, seed~42) for all Stage~2
proportions. Given the pilot scale of this experiment ($N{=}83$ total,
$N{=}30$ middle-position), we treat Stage~2 as a proof-of-concept
evaluation; exact significance tests are underpowered at this sample size and
are not reported.

%% ─────────────────────────────────────────────────────────────────────────────
\section{Results}
\label{sec:results}

\subsection{CLitM is Real and Large in Clinical EHR}

Across 2,196 instruction-response pairs and six model variants, retrieval
accuracy exhibits a pronounced U-shaped curve as a function of answer position
(Figure~\ref{fig:ucurve}). Peak accuracy of 59.5\% [95\% bootstrap CI: 52.4,
66.7] is observed at the 20--30\% position decile. Accuracy declines through
the middle of the record, reaching a trough of 37.6\% [31.2, 44.6] at the
70--80\% position decile---a \textbf{peak-to-trough gap of 21.9 percentage points}
(consistent with \citealt{liu2024lost} who report 15--25~pp using the same metric).
Average accuracy in middle-position instructions (30--70\% of EHR timeline) is
43.2\% versus 45.2\% in edge instructions---a 2.0~pp directional gap that
confirms systematic positional disadvantage for mid-timeline content even after
averaging out the outlier peak at 20--30\%.
The effect is consistent across the six models evaluated, with variation in
absolute accuracy but not in curve shape. Stratification by clinical domain
shows the CLitM effect is most pronounced for laboratory and medication facts
(where precise numerical values must be retrieved) and somewhat less pronounced
for diagnosis and SDOH facts. Examining the spatial distribution of reference
answers, 67.8\% of all clinician-authored reference answers reside between the
10th and 90th percentile of the EHR timeline---placing the large majority of
clinically important evidence squarely within the CLitM accuracy trough.

\begin{figure}[t]
  \centering
  \includegraphics[width=0.82\linewidth]{figures/exp1_litm_ucurve.png}
  \caption{Clinical LitM (CLitM) U-curve across 2,196 MedAlign
    instruction-response pairs and six language models. Peak accuracy 59.5\%
    [95\% CI: 52.4, 66.7] at 20--30\%; trough 37.6\% [31.2, 44.6] at
    70--80\%; gap = 21.9~pp. Shaded region indicates bootstrap 95\% CI.}
  \label{fig:ucurve}
\end{figure}

\subsection{Inhibitory Attention on Structured EHR Prediction}

Full results across four completed EHRSHOT tasks are reported in
Table~\ref{tab:ehrshot} (with approximate 95\% CIs via Hanley--McNeil). Differential
Transformer yields statistically significant improvements on the two most
class-imbalanced tasks: anemia (0.824 [0.818, 0.830] vs.\ 0.764 [0.757, 0.770],
$\Delta{=}{+}6.1$~pp AUROC; non-overlapping CIs) and hyperkalemia (0.615 [0.595, 0.635]
vs.\ 0.573 [0.554, 0.592], $\Delta{=}{+}4.2$~pp; borderline overlap).
On more balanced tasks, Differential Transformer is statistically significantly
\emph{below} standard attention: hyponatremia ($-1.9$~pp; CIs non-overlapping:
0.731 [0.726, 0.736] vs.\ 0.750 [0.745, 0.755]) and thrombocytopenia
($-1.5$~pp; non-overlapping: 0.865 [0.859, 0.870] vs.\ 0.880 [0.875, 0.885]).
H$_2$O substantially outperforms all models on hyperkalemia (0.737 [0.719, 0.756]),
the most severely class-imbalanced task, suggesting that KV cache compression
under sparse signal may inadvertently sharpen sensitivity to rare positive events.
The mean $\Delta$ Differential vs.\ Standard across four tasks is +1.7~pp AUROC
(range $-1.9$ to $+6.1$~pp). The task-specific pattern is consistent with the
hypothesis that inhibitory attention benefits sparse-signal tasks but may
modestly hurt balanced tasks where the second attention head suppresses
informative tokens.

\begin{table}[t]
\centering
\caption{EHRSHOT structured prediction results (test AUROC [95\% CI] / test AUPRC).
  Full dataset, 40 epochs, A10G GPU. Hypoglycemia excluded: H$_2$O run exceeded
  7,200~s wall-clock limit. CIs via Hanley--McNeil approximation (estimated test
  $n$ = 30\% of full dataset). \textbf{Bold} = best per task.}
\label{tab:ehrshot}
\scriptsize\setlength{\tabcolsep}{3.5pt}
\begin{tabular}{@{}p{2.3cm}p{3.6cm}p{3.6cm}p{3.2cm}@{}}
\toprule
Task & Standard Transformer & Differential Transformer & H$_2$O (keep 50\%) \\
\midrule
Anemia         & 0.764 [0.757, 0.770] / 0.814 & \textbf{0.824 [0.818, 0.830] / 0.852} & 0.763 [0.756, 0.770] / 0.824 \\
Hyperkalemia   & 0.573 [0.554, 0.592] / 0.033 & 0.615 [0.595, 0.635] / 0.083 & \textbf{0.737 [0.719, 0.756] / 0.124} \\
Hyponatremia   & \textbf{0.750 [0.745, 0.755]} / \textbf{0.557} & 0.731 [0.726, 0.736] / 0.539 & 0.734 [0.729, 0.739] / 0.536 \\
Thrombocytopenia & \textbf{0.880 [0.875, 0.885] / 0.815} & 0.865 [0.859, 0.870] / 0.797 & 0.866 [0.860, 0.871] / 0.798 \\
\midrule
\multicolumn{4}{l}{\small\textit{Mean $\Delta$ Diff vs.\ Std: +1.7~pp AUROC (range $-1.9$ to $+6.1$~pp)}} \\
\bottomrule
\end{tabular}
\end{table}

\subsection{Evidence Retrieval Baselines (Exp 3, Stage 1)}

Gold-answer retrieval results on the held-out test split (74 patients, 83 unique
instructions) are reported in Table~\ref{tab:qccsgate}. All three lexical and
semantic retrieval baselines substantially outperform the QCCS learned gate.
At $k{=}20$, BM25 achieves 98.8\% overall and 96.7\% in the middle
(30--70\% of EHR timeline) position band; dense retrieval achieves 94.0\%
overall and 90.0\% in the middle. The QCCS gate achieves 34.9\% overall and
23.3\% for middle-position instructions.

Critically, BM25 with section-header filtering achieves \emph{identical} recall
to standard BM25 (98.8\% overall, 96.7\% middle), confirming that the
section-header patterns are not a significant source of retrieval misses.
The QCCS gate underperforms all three retrieval baselines substantially.
Withholding the clinical question substantially reduces QCCS recall
(Appendix~\ref{app:noquery}), confirming that the gate prioritizes
query-relevant sentences rather than learning a generic text-quality score. Table~\ref{tab:ablation_k} shows that BM25 recall is largely
insensitive to $k$ above 10 ($\geq$96.7\% middle at all $k{\geq}10$), while
QCCS recall barely improves with $k$ in the middle band (23.3\% at all
$k{=}5$--$20$), suggesting that the gate's errors are systematic rather than
marginal.

The underperformance of the QCCS gate relative to retrieval baselines has a
structural explanation: the gate is trained with lexical overlap labels (sentences
sharing content words with the reference answer), which directly parallels the
BM25 term-frequency signal. The gate approximates---but does not exceed---the
lexical baseline at the retrieval stage. However, as Stage~2 shows, high
retrieval recall does not determine downstream LLM accuracy.

\begin{table}[t]
\centering
\caption{Gold-answer retrieval recall at $k{=}20$ on the held-out test split
  ($N{=}83$ unique instructions, 74 patients). Hit criterion: selected sentence
  contains at least one content word ($\geq$4 chars) from the reference answer.
  All retrieval methods include a fixed recency buffer of the 5 most-recent events.
  \textbf{Bold} = best.}
\label{tab:qccsgate}
\small
\begin{tabular}{@{}lccc@{}}
\toprule
Method & Overall & Middle (30--70\%) & Edge ($<$30\% or $>$70\%) \\
\midrule
BM25 (lexical) & \textbf{98.8\%} & 96.7\% & \textbf{100.0\%} \\
BM25-filtered (header-excl.) & \textbf{98.8\%} & 96.7\% & \textbf{100.0\%} \\
Dense (all-MiniLM-L6-v2) & 94.0\% & 90.0\% & 96.2\% \\
QCCS gate (learned) & 34.9\% & 23.3\% & 41.5\% \\
\bottomrule
\end{tabular}
\end{table}

\begin{table}[t]
\centering
\caption{$k$-sensitivity ablation: gold-answer retrieval recall (\%) vs.\ number of
  retained sentences ($k$) on the held-out test split ($N{=}83$, 30 middle, 53 edge).}
\label{tab:ablation_k}
\small
\setlength{\tabcolsep}{4pt}
\begin{tabular}{@{}rcccc ccccc@{}}
\toprule
 & \multicolumn{4}{c}{Overall recall (\%)} & & \multicolumn{4}{c}{Middle recall (\%)} \\
\cmidrule(lr){2-5}\cmidrule(lr){7-10}
$k$ & BM25 & BM25-filt & Dense & QCCS && BM25 & BM25-filt & Dense & QCCS \\
\midrule
 5 & 92.8 & 96.4 & 91.6 & 33.7 && 93.3 & \textbf{96.7} & 86.7 & 23.3 \\
10 & 97.6 & 97.6 & 94.0 & 33.7 && 96.7 & 96.7 & 90.0 & 23.3 \\
15 & 97.6 & 98.8 & 94.0 & 33.7 && 96.7 & 96.7 & 90.0 & 23.3 \\
20 & \textbf{98.8} & \textbf{98.8} & 94.0 & 34.9 && 96.7 & 96.7 & 90.0 & 23.3 \\
\bottomrule
\end{tabular}
\end{table}

% No-query ablation table moved to Appendix (see app:noquery)

\subsection{LLM Re-inference: Full Context vs.\ BM25 vs.\ QCCS}

End-to-end instruction-following results, evaluated by LLM-as-judge (primary
metric) and token-overlap (secondary cross-check), are reported in
Table~\ref{tab:llm}. All five arms are evaluated on the 83 unique (patient,
question) pairs from the held-out test split with Qwen2.5-7B-Instruct.
The full-context baseline uses a 16,384-token input window; compressing to
$k{=}20$ sentences reduces input length from $\sim$8k tokens (median) to
$\leq$500 tokens---a $\geq$16$\times$ reduction.

\textbf{QCCS is the only arm with consistently positive LLM accuracy.}
By LLM-as-judge scoring, QCCS achieves 20.0\% [95\% CI: 7, 37] for
middle-position instructions (30--70\% of EHR timeline) versus 0.0\% [0, 0]
for dense retrieval, 0.0\% [0, 0] for BM25-filtered, 3.3\% [0, 10] for BM25,
and 6.7\% [0, 17] for full context. Overall, QCCS reaches 26.5\% [17, 36]
versus $\leq$3.6\% for all comparators. The pattern is consistent across
position bands: QCCS is the best-performing arm in every decile
(Table~\ref{tab:llm}), with the largest margin at early positions
(0--10\%: QCCS 46.2\% vs.\ all comparators 0.0\%).

\textbf{Dense retrieval and BM25 both fail at Stage~2.}
Dense retrieval (all-MiniLM-L6-v2) achieves 90.0\% Stage~1 recall in the
middle band yet 0.0\% LLM accuracy there---identical to BM25-filtered and
worse than standard BM25 (3.3\%). This result is striking: a semantically rich
retriever with near-BM25 recall provides no benefit over full-context
inference for downstream accuracy. The full-context baseline scores near-zero
across most position bands, with Qwen2.5-7B-Instruct failing to retrieve
evidence at essentially any position despite it being present in the 16k-token
window, making context compression a practical necessity.

\textbf{Token-overlap vs.\ semantic scoring.} Token-overlap inflates dense
and BM25-filtered in the 30--50\% band (33.3\% each by token-overlap vs.\ 0.0\%
by judge) and the baseline (11.1\% vs.\ 11.1\% aligned here, but 0\%
judge-vs-tok discrepancies elsewhere). LLM-as-judge scoring removes these
false positives and confirms that future EHR context selection evaluations
should use semantic metrics.

\textbf{BM25 and dense failure analysis.} BM25 selects EHR sentences by
lexical overlap with the query; many MedAlign EHR records contain structured
``Question:''/``Answer:'' formatted sections (social history, review of
systems) whose headers match query tokens. BM25-filtered removes the most
obvious header patterns but achieves identical Stage~1 recall (98.8\%) and
0.0\% middle LLM accuracy, showing that the failure is not addressable by
header filtering. Dense retrieval avoids these lexical traps yet still fails,
suggesting the problem is not lexical false positives per se but rather that
both methods retrieve \emph{contextually plausible} but not
\emph{query-conditionally relevant} sentences. QCCS, trained on
query-conditioned lexical overlap labels, learns a feature space that
distinguishes query-relevant from coincidentally relevant context, providing
semantically coherent, question-specific selection. This confirms the central
claim: \emph{retrieval recall does not determine downstream LLM accuracy;
context precision does.}

Given the pilot scale of this experiment ($N{=}83$ total, $N{=}30$
middle-position), differences are not statistically significant
(bootstrap CIs overlap substantially). These findings should be interpreted
as proof-of-concept estimates requiring confirmation in a larger evaluation.

\begin{table}[t]
\centering
\caption{Qwen2.5-7B-Instruct end-to-end accuracy (\%) on MedAlign held-out test split
  ($N{=}83$ unique instructions, 74 patients; 16k-token context window).
  \textbf{Primary metric: LLM-as-judge (Claude Haiku semantic evaluation).}
  Secondary: token-overlap ($\geq$50\% expected tokens present).
  All five arms use identical inference code; compressed arms (BM25, BM25f,
  Dense, QCCS) use top-20 sentences plus a last-5 recency buffer.
  BM25f = BM25 with section-header filtering. 95\% bootstrap CIs for summary rows
  (5,000 resamples, seed~42). \textbf{Bold} = best per row (judge metric).}
\label{tab:llm}
\small
\setlength{\tabcolsep}{2pt}
\begin{tabular}{@{}lc ccccc ccccc@{}}
\toprule
 & & \multicolumn{5}{c}{LLM-as-judge (\%, primary)} & \multicolumn{5}{c}{Tok-overlap (\%, secondary)} \\
\cmidrule(lr){3-7}\cmidrule(lr){8-12}
Pos.\ band & $N$ & Full & BM25 & B25f & Dns & QCCS & Full & BM25 & B25f & Dns & QCCS \\
\midrule
0--10\%   & 13 &  0.0 &  0.0 &  0.0 &  0.0 & \textbf{46.2} &  0.0 &  0.0 &  0.0 &  0.0 & 15.4 \\
10--30\%  & 12 &  0.0 &  0.0 &  0.0 &  0.0 & \textbf{25.0} &  0.0 &  0.0 &  0.0 &  0.0 &  8.3 \\
30--50\%  &  9 & 11.1 &  0.0 &  0.0 &  0.0 & \textbf{22.2} & 11.1 &  0.0 & 22.2 & 33.3 & 33.3 \\
50--70\%  & 21 &  4.8 &\textbf{4.8}&  0.0 &  0.0 & \textbf{19.0} &  0.0 &  4.8 &  4.8 &  0.0 &  9.5 \\
70--90\%  & 13 &  0.0 &  0.0 &  7.7 &  0.0 & \textbf{30.8} &  7.7 &  0.0 &  0.0 &  7.7 &  7.7 \\
90--100\% & 15 &  6.7 &  0.0 &  0.0 & 13.3 & \textbf{20.0} &  6.7 &  0.0 &  0.0 &  0.0 &  6.7 \\
\midrule
Middle (30--70\%) & 30 & 6.7\,[0,17] & 3.3\,[0,10] & 0.0\,[0,0] & 0.0\,[0,0] & \textbf{20.0\,[7,37]} &  3.3 &  3.3 & 10.0 & 10.0 & 16.7 \\
Edge              & 53 & 1.9\,[0,6]  & 0.0\,[0,0]  & 1.9\,[0,6] & 3.8\,[0,9] & \textbf{30.2\,[19,43]}&  3.8 &  0.0 &  0.0 &  1.9 &  9.4 \\
Overall           & 83 & 3.6\,[0,8]  & 1.2\,[0,4]  & 1.2\,[0,4] & 2.4\,[0,6] & \textbf{26.5\,[17,36]}&  3.6 &  1.2 &  3.6 &  4.8 & 12.0 \\
\bottomrule
\end{tabular}
\end{table}

\begin{figure}[t]
  \centering
  \includegraphics[width=0.82\linewidth]{figures/exp3_llm_ucurve.png}
  \caption{End-to-end LLM instruction-following accuracy (\%, LLM-as-judge primary metric)
    by answer position decile: full context (16k-token window), BM25 top-20,
    and QCCS top-20 (Qwen2.5-7B-Instruct, $N{=}83$ unique instructions; three-arm
    subset of five-arm experiment; see Table~\ref{tab:llm} for full five-arm results).
    QCCS is the best-performing arm in every position band; overall judge accuracy
    26.5\% [17, 36] vs.\ $\leq$3.6\% for all comparators.
    Neither BM25 (96.7\% Stage~1 recall) nor dense retrieval (90.0\% recall) translates
    to LLM accuracy, confirming that context precision---not retrieval recall---is the
    binding constraint.}
  \label{fig:llm}
\end{figure}

%% ─────────────────────────────────────────────────────────────────────────────
\section{Discussion}
\label{sec:discussion}

The 21.9 percentage-point CLitM gap between peak and trough positions is not
a benchmark artifact: 67.8\% of reference answers reside between the 10th and
90th EHR percentile, so this gap maps directly onto content clinicians consider
important.

\paragraph{Retrieval recall does not predict LLM accuracy.}
At the retrieval level, BM25 and dense retrieval substantially outperform QCCS
(98.8\%, 94.0\% vs.\ 34.9\% overall; Table~\ref{tab:qccsgate}). At the
downstream LLM accuracy level, the ranking inverts: QCCS 26.5\% [17, 36]
vs.\ $\leq$3.6\% for all comparators (Table~\ref{tab:llm}). This reversal
establishes that context selection systems must be evaluated end-to-end with
semantic metrics, not retrieval recall alone.

\paragraph{QCCS provides consistent improvement across position bands.}
QCCS outperforms all comparators in \emph{every} position decile, including
early-context (0--10\%: 46.2\% vs.\ 0.0\% for all others), where primacy bias
might have favored the full-context baseline. The full-context Qwen2.5-7B
fails even at edge positions, confirming that 7B-scale models lack sufficient
retrieval capacity to benefit from positional advantages in long EHRs.

\paragraph{Structured prediction: selective benefit of inhibitory attention.}
Differential Transformer yields significant gains on class-imbalanced tasks
(anemia: +6.1~pp AUROC; hyperkalemia: +4.2~pp), but is significantly
\emph{worse} on more balanced tasks (hyponatremia: $-1.9$~pp;
thrombocytopenia: $-1.5$~pp; non-overlapping CIs in both cases). The benefit
is real but selective, arising when target signal is sparse relative to background.

\paragraph{Limitations.}
See Appendix~\ref{app:limitations} for a full discussion of limitations including
single-site data, pilot-scale Stage~2 evaluation, QCCS training circularity, and
context coverage confounds.

%% ─────────────────────────────────────────────────────────────────────────────
\section{Conclusion}
\label{sec:conclusion}

CLitM is a robust clinical risk (21.9~pp gap; 67.8\% of reference evidence at
mid-timeline positions). QCCS achieves 7$\times$ higher end-to-end LLM accuracy
than BM25/dense retrieval despite lower Stage~1 recall, establishing context
precision as the binding constraint for EHR context selection.

%% ─────────────────────────────────────────────────────────────────────────────
\bibliographystyle{plainnat}
\bibliography{paper}

%% ─────────────────────────────────────────────────────────────────────────────
%% NEURIPS PAPER CHECKLIST
%% ─────────────────────────────────────────────────────────────────────────────
\newpage
\section*{NeurIPS Paper Checklist}

\begin{enumerate}

\item {\bf Claims}
  \item[] Question: Do the main claims made in the abstract and introduction
    accurately reflect the paper's contributions and scope?
  \item[] Answer: \answerYes{}
  \item[] Justification: All quantitative claims in the abstract (21.9~pp CLitM
    gap; BM25 98.8\% overall vs.\ QCCS 34.9\% at $k{=}20$; LLM-as-judge
    five-arm comparison showing QCCS 26.5\% [17,36] vs.\ $\leq$3.6\% for all
    comparators; +6.1~pp anemia AUROC with Differential Transformer) are
    directly reported in Tables~\ref{tab:qccsgate}--\ref{tab:llm} and
    Table~\ref{tab:ehrshot}, backed by the experiments in Section~3. Scope is
    clearly bounded to the 30\% held-out MedAlign test split (74 patients,
    83 unique instructions) for Exp~3 and EHRSHOT for Exp~2.
  \item[] Guidelines: N/A

\item {\bf Limitations}
  \item[] Question: Does the paper discuss the limitations of the work?
  \item[] Answer: \answerYes{}
  \item[] Justification: Limitations are discussed explicitly in
    Section~5 (Discussion, Limitations paragraph): small MedAlign sample
    ($n{=}275$ patients), single academic medical center population, proof-of-concept
    Stage~2 sample size ($N{=}83$, wide CIs), QCCS training criterion circularity,
    16k-token context truncation, and the need for prospective clinical validation.

\item {\bf Theory Assumptions and Proofs}
  \item[] Question: For each theoretical result, does the paper provide the
    full set of assumptions and a complete proof?
  \item[] Answer: \answerNA{}
  \item[] Justification: This is an empirical paper. The QCCS architecture is
    presented as a formal model in Equations~\ref{eq:gate}--\ref{eq:qccs};
    no theoretical proofs are claimed.

\item {\bf Experimental Result Reproducibility}
  \item[] Question: Does the paper fully disclose all the information needed
    to reproduce the main experimental results of the paper to the extent that
    it affects the main claims?
  \item[] Answer: \answerYes{}
  \item[] Justification: Section~3 provides full experimental details including
    hyperparameters, hardware, train/test splits, evaluation metrics, and scoring
    criteria. Code is released at
    \texttt{[anonymized for review]}; data access
    instructions are in the repository.

\item {\bf Open access to data and code}
  \item[] Question: Does the paper provide open access to the data and code
    used in the paper?
  \item[] Answer: \answerYes{}
  \item[] Justification: Experiment code is publicly available at
    \texttt{[anonymized for review]} under the MIT
    License. MedAlign and EHRSHOT are available under their respective data use
    agreements (see repository \texttt{data/README.md} for access instructions).
    Protected health information is not released.

\item {\bf Experimental Setting/Details}
  \item[] Question: Does the paper specify all the training and test
    details (e.g.\ data splits, hyperparameters, how they were selected) to
    support the claims?
  \item[] Answer: \answerYes{}
  \item[] Justification: All details provided: A10G GPU hardware, 40 epochs
    for Exp~2, AdamW lr~$= 10^{-3}$ batch~32 epochs~15 for QCCS gate,
    70/30 patient-level split for Exp~3, $k=20$ top-sentence retention;
    Stage~2 primary metric is LLM-as-judge (Claude Haiku), secondary is
    $\geq$50\% token-overlap; bootstrap 95\% CIs reported.

\item {\bf Experiment Statistical Significance}
  \item[] Question: Does the paper report error bars suitably and correctly
    defined?
  \item[] Answer: \answerYes{}
  \item[] Justification: Bootstrap 95\% CIs (5,000 resamples) are reported for
    all Stage~2 LLM accuracy estimates (Table~\ref{tab:llm}, summary rows).
    Hanley--McNeil 95\% CIs are reported for AUROC in Table~\ref{tab:ehrshot}.
    Bootstrap 95\% CIs are shown in Figure~\ref{fig:ucurve} for Experiment~1.
    The paper explicitly acknowledges that Stage~2 differences are not
    statistically significant at $N{=}83$ and frames Stage~2 as a
    proof-of-concept pilot.

\item {\bf Experiments Compute Resources}
  \item[] Question: For each experiment, does the paper provide sufficient
    information on the computational resources (type of compute, compute time)
    used?
  \item[] Answer: \answerYes{}
  \item[] Justification: Experiments run on Modal GPU instances (Exp~2: A10G
    24~GB; Exp~3 LLM inference: A100 80~GB). Exp~2 total compute $\sim$4--8
    GPU-hours per task. Exp~3 gate trains on CPU in $\sim$15 minutes;
    LLM re-inference requires $\sim$4--6 GPU-hours on A100 (three-arm,
    $N{=}83$).

\item {\bf Code of ethics}
  \item[] Question: Does the research conducted in the paper conform, in every
    respect, with the NeurIPS Code of Ethics?
  \item[] Answer: \answerYes{}
  \item[] Justification: No human subjects were recruited for this study.
    All patient data used is de-identified and accessed under data use
    agreements. Appendix~\ref{sec:broader} explicitly addresses potential
    clinical harms and the need for human oversight.

\item {\bf Broader impacts}
  \item[] Question: Does the paper discuss both potential positive societal
    impacts and negative societal impacts of the work performed?
  \item[] Answer: \answerYes{}
  \item[] Justification: Appendix~\ref{sec:broader} (Broader Impact) discusses
    positive impacts (improving reliability of clinical AI, reducing diagnostic
    errors) and negative impacts (potential for adversarial exploitation of
    attention biases).

\item {\bf Safeguards}
  \item[] Question: Does the paper describe safeguards that have been put in
    place for responsible release of data or models?
  \item[] Answer: \answerYes{}
  \item[] Justification: No raw EHR data is released. Code does not contain
    patient data. Data access requires signing data use agreements.
    The paper notes human oversight requirements for any deployment.

\item {\bf Licenses for existing assets}
  \item[] Question: Are the creators or original owners of assets used
    credited and are the license terms met?
  \item[] Answer: \answerYes{}
  \item[] Justification: MedAlign (Stanford AIMI, DUA), EHRSHOT (PhysioNet
    credentialing), and all software dependencies are credited in the text
    and bibliography.

\item {\bf New Assets}
  \item[] Question: Are new assets introduced in the paper well documented
    and is the documentation provided alongside the assets?
  \item[] Answer: \answerYes{}
  \item[] Justification: The QCCS gate implementation is documented and
    released at the GitHub repository with a README, requirements, and
    data access instructions.

\item {\bf Crowdsourcing and research with human subjects}
  \item[] Question: For crowdsourcing experiments and research with human
    subjects, does the paper include the full text of instructions given to
    participants and screenshots, if applicable, as well as details about
    compensation (if any)?
  \item[] Answer: \answerNA{}
  \item[] Justification: No crowdsourcing or research with human subjects was
    conducted. This work uses existing annotated datasets (MedAlign, EHRSHOT).

\item {\bf Institutional review board (IRB) approvals or equivalent for research with human subjects}
  \item[] Question: Does the paper describe potential risks incurred by study
    participants, whether such risks were disclosed to the subjects, and whether
    IRB approvals (or an equivalent approval/review) were obtained?
  \item[] Answer: \answerNA{}
  \item[] Justification: No human subjects research was conducted. All EHR
    data is de-identified and accessed under existing data use agreements
    (MedAlign: Stanford AIMI DUA; EHRSHOT: PhysioNet credentialing). No new
    IRB approval was required for this secondary analysis.

\item {\bf Declaration of LLM usage}
  \item[] Question: Does the paper describe the usage of LLMs if it is an
    important, original, or non-standard component of the core methods in
    this research?
  \item[] Answer: \answerYes{}
  \item[] Justification: LLMs are used in two core components: (1)
    Qwen2.5-7B-Instruct for Stage~2 re-inference evaluation (Section~3.3);
    (2) Claude Haiku (\texttt{claude-haiku-4-5-20251001}) as LLM-as-judge
    for semantic evaluation of response quality (Section~3.3). Both usages
    are described in detail in the Methods section.

\end{enumerate}

%% ─────────────────────────────────────────────────────────────────────────────
\appendix

\section{Broader Impact}
\label{sec:broader}

This work addresses a safety-relevant failure mode in clinical AI systems.
Uncorrected positional bias in clinical decision support tools could contribute
to diagnostic errors and medication-related adverse events. QCCS reduces this
risk by directing model attention toward query-relevant evidence regardless of
temporal position. The absolute accuracy levels demonstrated (up to 46.2\% with
QCCS for early-context positions) are insufficient for unsupervised clinical
decision-making; human oversight remains essential. Positional bias mitigation
is a necessary but not sufficient condition for safe clinical AI deployment.
There are potential negative uses of techniques for identifying attention biases
(e.g., deliberately placing adversarial information in the middle of a context
to be missed). Transparent characterization and architectural mitigation of this
effect reduces net risk by enabling developers to test for and address positional
bias in production systems.

\section{Limitations}
\label{app:limitations}

MedAlign contains 275 patients from a single academic medical center, limiting
statistical power and generalizability. The Stage~2 LLM evaluation ($N{=}83$,
$N{=}30$ middle) is a proof-of-concept pilot; differences between arms are not
statistically significant at this sample size. The QCCS gate training criterion
(lexical overlap) creates a structural circularity: the gate learns to approximate
BM25 relevance, which may overstate the difficulty of semantic selection under
richer supervision (e.g., LLM-generated labels). The 16,384-token full-context
baseline drops late EHR events on long records; BM25 and QCCS select from the full
stream, so part of their advantage reflects coverage rather than CLitM mitigation.
Prospective clinical validation across diverse populations remains necessary.

\section{Extended CLitM Curves by Model and Clinical Domain}
\label{app:ucurves}

Figure~\ref{fig:bymodel} shows the CLitM U-curve stratified by each of the six
language models evaluated in Experiment~1. The qualitative U-shaped pattern is
consistent across all models. GPT-4-class models achieve higher absolute
accuracy (peak $\sim$65--70\%) than open-weight 7B--13B models (peak
$\sim$45--55\%), but both exhibit a pronounced trough at the 70--80\% position
decile. Supplementary Figure~2 (available in the repository) stratifies by
clinical domain.

\begin{figure}[h]
  \centering
  \includegraphics[width=\linewidth]{figures/exp1_litm_by_model.png}
  \caption{CLitM U-curves stratified by language model. All six models
    exhibit the same qualitative U-shaped positional bias.}
  \label{fig:bymodel}
\end{figure}

\section{QCCS Gate Visualization Examples}
\label{app:gate}

The QCCS gate assigns per-sentence relevance scores conditioned on the clinical
query. High-scoring sentences (retained in the compressed context) typically
correspond to sections of the EHR record explicitly named in the instruction
(e.g., medication lists for medication queries; laboratory sections for lab
queries). Gate visualization examples are available in the code repository
at \texttt{[anonymized for review]}.

\section{Full EHRSHOT Results Including Validation AUROC}
\label{app:ehrshot}

Validation AUROC values for Table~\ref{tab:ehrshot} (reported for reference;
test AUROC and AUPRC are the primary metrics):

\begin{center}
\small
\begin{tabular}{@{}lccc@{}}
\toprule
Task & Standard (val AUROC) & Differential (val AUROC) & H$_2$O (val AUROC) \\
\midrule
Anemia         & 0.878 & 0.888 & 0.878 \\
Hyperkalemia   & 0.607 & 0.654 & 0.776 \\
Hyponatremia   & 0.747 & 0.732 & 0.736 \\
Thrombocytopenia & 0.879 & 0.863 & 0.863 \\
\bottomrule
\end{tabular}
\end{center}

\section{Query Conditioning Ablation}
\label{app:noquery}

Table~\ref{tab:ablation_noq} reports the QCCS gate retrieval recall with and
without query conditioning, evaluated on the full MedAlign test population
($N{=}690$ TSV rows across 74 test patients before instruction deduplication).
Query conditioning contributes +14.7~pp overall and +23.8~pp for middle-position
recall, confirming that query context provides a meaningful training signal and
that the gate is not simply learning a generic text-quality score. Note that
the absolute values (41.7\%/38.1\% with query) are higher than those reported
in the main paper's Table~\ref{tab:qccsgate} (34.9\%/23.3\%) because the
ablation uses the pre-deduplication test population; the direction and magnitude
of the query conditioning effect are robust across both analyses.

\begin{table}[H]
\centering
\caption{No-query ablation at $k{=}20$: QCCS gate with and without query
  conditioning (full test population, $N{=}690$ rows before deduplication).}
\label{tab:ablation_noq}
\small
\begin{tabular}{@{}lcc@{}}
\toprule
Condition & Overall & Middle (30--70\%) \\
\midrule
QCCS (with query) & 41.7\% & 38.1\% \\
QCCS (no query)   & 27.0\% & 14.3\% \\
\midrule
Query contribution & +14.7~pp & +23.8~pp \\
\bottomrule
\end{tabular}
\end{table}

\end{document}
